\section[Numérisation à Cujas]{Les balbutiements de la numérisation à Cujas (2001 – 2007)}

Le choix de la bibliothèque Cujas est d'orienter la valorisation de son contenu vers \enquote{un lectorat universitaire de haut niveau}. Noëlle Balley et Sébastien Dalmon expliquent la raison derrière cette volonté venant de la spécialisation du fonds. Pour eux, le contenu est \enquote{souvent aride} et donc peu propice à une \enquote{valorisation orientée vers le grand public}. La bibliothèque Cujas privilégie donc des corpus pour le lectorat unversitaire de haut niveau, en mettant à leur disposition des \enquote{outils facilitant leur travail sur les textes}. Toutefois, ils soulignent que malgré les moyens mis en oeuvre pour favoriser un lectorat élitiste, pour autant il ne faut pas \enquote{négliger le grand public}.

L'un des premiers projets de numérisation de la bibliothèque Cujas fut dans les années 2000 et 2001. A cette époque, Cujas numérise 17 titres représentant environ 22 volumes. Elle le fait réaliser par la société Safig. La bibliothèque demande au prestataire de lui numériser ses documents \enquote{en mode image, avec une navigation grâce à la pagination}. La société Safig offre à Cujas la possibilité d'héberger ses numérisations son site. La bibliothèque laisse donc ses numérisations chez le prestataire qui s'occupait alors de diffuser les documents en ligne. Malheureusement, la société est victime d'une attaque des serveurs, entraînant une perte considérable des documents numérisés. La bibliothèque tente de récupérer ce qu'elle peut. Safig envoie ce qui lui reste \enquote{sous forme de CD-ROM}. A la réception des CD-ROM, le personnel de Cujas tente de voir s'il est possible \enquote{de réutiliser du contenu}. Hélas, le contenu est irrécupérable, voir pour certain corrompu, la décision est prise finalement de tout renumériser.

Le 5 et 6 novembre 2008, la bibliothèque Cujas avait organisé un \enquote{colloque international en collaboration avec le Sénat},portant sur le doyen Jean Carbonnier. Noëlle Balley et Sébastien Dalmon parlent de ce projet comme \enquote{la première entreprise de numérisation effectivate a concerné, pour des raisons conjoncturelles, les polycopiés du doyen Jean Carbonnier}. La numérisation de ses polycopiés est réalisé à \enquote{l'occasion du centenaire de la naissance de ce grand juriste civiliste}. A l'origine, ce projet réside dans une \enquote{participation modeste} de la bibliothèque Cujas à \enquote{un projet éditorial}. L'anthropologue du droit Raymond Verdier initie ce projet. Son objectif est \enquote{de rassembler les écrits de Jean Carbonnier introuvables, inédits ou dispersés}. La valorisation de ce projet prend différentes formes : les colloques qui ont lieu au Sénat, des actions de commmunication, une exposition virtuelle.

En parallèle des colloques, la bibliothèque Cujas organise une exposition virtuelle. Elle a décidé \enquote{de mettre en ligne une exposition virtuelle consacrée à Jean Carbonnier, constatant l'absence, sur Internet, de ressources suffisamment fiables et complètes sur le Doyen}. Leur but est de faire connaitre \enquote{une grande figure du droit contemporain, méconnue du grand public, mais aussi à répondre aux attentes des spécialistes}. Cujas a joué sur son expérience des expositions virtuelles pour réaliser celle-ci afin de mettre en avant la figure de Jean Carbonnier. En effet, c'est loin d'être sa première exposition virtuelle, la bibliothèque en avait déjà réalisé d'autres avant celle-ci. Elle avait réalisé \enquote{celle sur le Bicentenaire du Code civil (1804-2004) en collaboration avec la Cour de cassation et l'Ordre des avocats}.

L'exposition virtuelle offre un \enquote{accès à treize polycopiés de cours du doyen Carbonnier professés entre 1958 et 1974, portant sur la sociologie juridique ou le droit civil}. Ce corpus de documents peut être retrouvé sur la bibliothèque numérique actuelle, anciennement CujasNum, intitulé \enquote{Cours de droit Carbonnier}. La mission de numérisation de la bibliothèque fut un énorme travail. En effet, ils ont numérisés \enquote{en mode texte, avec une relectur intégrale}. Cette numérisation intense a ammené la création d'une \enquote{plate-forme de recherche [...]par l'ingénieur informatique de la bibliothèque; des sommaires dynamiques ont également été créés}. A cette époque, la bibliothèque Cujas a à coeur de donner le  \enquote{meilleur accès aux informations contenues dans le texte} qu'elle s'occupe elle-même de la correction intégrale des Optical Recognition Characters, corrigeant à la main les erreurs afin d'obtenir un taux de 100\%. Le fonds Carbonnier est la première collection ajoutée dans la bibliothèque numérique. Dans un souci d'apporter un confort aux chercheurs, l'interface de recherche a été amélioré. Elle permet de croiser les \enquote{recherches sur les métadonnées, la table des matières et le plein texte, soit en favorisant la sérendipidité grâce à une navigation par listes d'auteurs ou de titres, ou encore à travers un thésaurus spécialement conçu pour la bibliothèque numérique}. A ce moment-là, la bibliothèque est fière car CujasNum répond aux mieux aux besoins de ces chercheurs, plus tard elle ajoutera une visionneuse permettant l'accès \enquote{à une page précise sans télécharger tout le document}.

La numérisation de la collection Jean Carbonnier est la première opération de numérisation de la bibliothèque Cujas. Elle a donc permis de tester \enquote{en grandeur réelle, la chaîne de numérisation qu'elle mettait en place pour son propre programme}. En effet, celle-ci sera réemployée au moment de la numérisation de son deuxième plus gros projet de numérisation, la liste Pfister-Roumy.